\section{\ref{sec:gaba}장. \nt{GLUT}와 \nt{GABA}의 협업}

이 장의 논리적, 동역학적 명제는 선형 해석과 옴의 법칙으로 장 안에서 유도했다. 실험은 이 명제들이 기대는 회로, 즉 지연된 금지, 흥분 뒤에 따라오는 직접 억제, 억제에 의한 안정화, \nt{GLUT}--\nt{GABA} 루프의 리듬이 뇌에서 실제로 작동한다는 것을 보여 준다.

{\footnotesize
\setlength{\tabcolsep}{3pt}\renewcommand{\arraystretch}{1.15}
\begin{longtable}{>{\raggedright}p{0.5cm}>{\raggedright}p{4.0cm}>{\raggedright}p{5.6cm}>{\raggedright}p{2.3cm}>{\raggedright\arraybackslash}p{1.7cm}}
\toprule
No. & 명제 & 실험과 측정한 것 & 출처 & 상태 \\
\midrule\endfirsthead
\toprule
No. & 명제 & 실험과 측정한 것 & 출처 & 상태 \\
\midrule\endhead
1 & 피질 뉴런의 5분의 1에서 4분의 1이 \nt{GABA} 뉴런이다 & 원숭이 피질 $10$개 영역의 GABA 면역염색: $8$개 영역에서 뉴런의 약 $25\,\%$, 일차 시각 피질에서는 $20\,\%$ 미만. 피질 억제성 뉴런의 다양성에 관한 총설 & \cite{Hendry1987,Markram2004} & 측정됨 \\
2 & 억제가 하는 일은 붙는 자리가 크게 좌우한다. 가지돌기, 세포체, 스파이크가 생기는 자리, 다른 뉴런 전선의 말단 & 총설: 피질 \nt{GABA} 뉴런의 부류는 수신자 위의 표적 자리로 구별된다. 피라미드 뉴런의 세포체와 가지돌기에서 동시 기록: 직접 억제는 가지돌기보다 세포체에서 훨씬 강하다. 척수에서 다른 뉴런 전선의 말단에 가해지는 억제는 입력 선로 하나의 신경전달물질 방출을 줄인다(측정 총설) & \cite{Markram2004,Pouille2001,Rudomin1999} & 측정됨 \\
3 & 중개자를 거친 부호 전환의 값은 지연 하나, 약 1밀리초이다. 흥분 뒤에 따라오는 억제는 일치 창을 좁힌다 & 쥐 해마 절편의 피라미드 뉴런: 중개자 하나를 거친 억제가 직접 흥분 뒤에 짧은 지연으로 따라오며, 두 흥분성 입력이 문턱값까지 합산되는 창은 $2\ms$보다 짧다. 이 억제가 약한 가지돌기에서는 창이 넓다 & \cite{Pouille2001} & 측정됨 \\
4 & 금지 $a\wedge\bar b$~\eqref{eq:veto}는 OR, 상수 1과 함께 함수적으로 완전하다(명제~\ref{prop:gabacomplete}) & 장에서의 증명. 고전적 결과: 억제성 입력이 있는 문턱 소자의 네트워크는 어떤 논리식이든 구현한다. 모델에 관한 명제이며, 실험은 없다 & \cite{McCullochPitts1943} & 이론 \\
5 & 지연된 금지는 최적 속도 $v^*=\Delta x/d$~\eqref{eq:vstar}의 운동 방향 검출기를 준다 & 토끼 망막: 방향 선택성은 “무효” 방향으로 움직일 때 흥분을 끄는 억제로 만들어진다. 선택성 세포의 흥분성 입력과 억제성 입력을 따로 기록한 결과, 별 모양 아마크린 세포(\nt{GABA})가 “무효” 쪽을 향한 돌기로만 비대칭적으로 억제한다는 것이 드러났다. $v^*$ 공식은 장에서 유도했다 & \cite{Barlow1965,Fried2002} & 측정됨 \\
6 & 분로 억제는 전압에서 빼지 않고 전압을 나눈다~\eqref{eq:shunt}(그림~\ref{fig:shunt}) & 염소 평형 전위가 휴지 전위 근처일 때 세 컨덕턴스로 된 분배기에 대한 옴의 법칙. 스파이크를 내지 않는 뉴런에 관한 것이다 & 이 장의 유도 & 이론 \\
7 & 분로는 스파이크 발화율에 뺄셈으로 작용하고, 나눗셈은 억제와 균형 잡힌 배경 잡음의 조합에서 나온다 & 모델: 발화하는 뉴런에서는 분로를 지나는 전류가 거의 일정하여 발화율에 대한 효과가 뺄셈이다. 실험: 쥐 체성감각 피질(절편)의 피라미드 뉴런에 동적 클램프로 흥분성, 억제성 배경 컨덕턴스를 주입했다. 둘을 동시에 균형 있게 늘리면 발화율의 뚜렷한 이동 없이 응답 기울기가 나누어진다. 총설: 총활동으로 나누는 연산은 뇌의 여러 부위에서 나타난다 & \cite{HoltKoch1997,Chance2002,Carandini2012} & 측정됨 \\
8 & $\beta>1$이면 상호 억제가 승자 하나를 고른다(명제~\ref{prop:wta}, \eqref{eq:wta}) & 고윳값 $-(1\pm\beta)/\tau$는 장에서 유도했다. $\beta=0.5$에서 발화율 $0.73$과 $0.53$은 고정점 $r_{1,2}=(I_{1,2}-\beta I_{2,1})/(1-\beta^2)$이다. \texttt{models/}에 스크립트는 없다. 장에서 직접 실험을 제시하지 않는다 & 이 장의 유도 & 이론 \\
9 & \nt{GABA}를 거친 신호로 기억을 초기화; 상호 억제하는 두 무리는 래치이다 & 그림~\ref{fig:bistable}과 명제~\ref{prop:wta}에서 나온다. 실험은 없다 & 이 장의 유도 & 이론 \\
10 & \nt{GLUT}--\nt{GABA} 루프의 평형은 억제가 충분히 강하고 충분히 빠르면 안정하다(명제~\ref{prop:ei}) & 두 집단 모델~\eqref{eq:ei}. 조건 (i)과 (ii)는 그 행렬의 행렬식과 대각합이며 장에서 유도했다 & \cite{WilsonCowan1972} & 이론 \\
11 & $w_{EE}=1.5$, $w_{EI}w_{IE}=2$에서 빠른 억제($\tau_I=2\ms$)는 $E^*=0.67h$를 유지하고, 느린 억제($30\ms$)는 진동을 부추긴다(그림~\ref{fig:ei}) & \eqref{eq:ei}의 수치해, 스크립트 \texttt{figures/make\_ei.py} & 계산 & 모델 \\
12 & 피질은 억제 안정화 체제에서 작동한다. 특징은 \nt{GABA} 뉴런을 밀어 주면 그 발화율이 떨어진다는 것이다~\eqref{eq:isn} & 이론이 역설적 응답을 예측했다. 실험: 고양이 일차 시각 피질의 세포내 기록에서 주변 자극으로 응답이 억눌릴 때 억제성 컨덕턴스는 늘지 않고 흥분성 컨덕턴스와 함께 줄어든다. 생쥐의 시각, 체성감각, 운동 피질에서 PV 뉴런을 광유전학으로 흥분시키면 자극이 없을 때와 가벼운 마취에서도 억제성 뉴런의 발화율이 낮아진다. 그림~\ref{fig:ei}의 수치에서 계수 $-1/3$은 장에서 유도했다 & \cite{Tsodyks1997isn,Ozeki2009,Sanzeni2020} & 측정됨 \\
13 & “\nt{GLUT}가 \nt{GABA}를 흥분시키고, \nt{GABA}가 \nt{GLUT}를 억제한다”는 루프는 주기 $12$--$50\ms$의 빠른 리듬을 낳는다 & 생쥐 체성감각 피질에서 빠른 \nt{GABA} 뉴런을 $8$--$200$\,Hz로 광유전학 자극하면 감마 리듬($20$--$80$\,Hz, 주기 $12$--$50\ms$)이 선택적으로 강해진다. \nt{GLUT} 뉴런 자극은 더 느린 리듬만 강화한다 & \cite{Cardin2009,Buzsaki2012} & 측정됨 \\
\bottomrule
\end{longtable}}
